\begin{name}
	{ÔN TẬP KIỂM TRA CUỐI HỌC KÌ 1}
	{TOÁN 10}
	{LỚP TOÁN THẦY PHÁT}
	{\thoigian}
\end{name}

%================================PHẦN I
\caulc
\Opensolutionfile{ans}[ans/0-HK1-CD-1-SoBacNinh-2324-Phan-I]
%Thay Khối lớp thành 10/11, Đề mẫu thành thứ tự đề

%%%============EX_1==============%%%
\begin{ex}%[0-HK1-CD-1-SoBacNinh-2324]%[VN-MT-6, Tống Văn Ký]%[0D3H1-2]
	Tập xác định của hàm số $y=\dfrac{2}{\sqrt{x+1}}$ là
	\choice
	{$\mathbb{R}\setminus \{-1\}$}
	{$[-1;+\infty)$}
	{\True $(-1;+\infty)$}
	{$(1;+\infty)$}
	\loigiai{
		Hàm số $y=\dfrac{2}{\sqrt{x+1}}$ xác định $\Leftrightarrow x+ 1 > 0 \Leftrightarrow x > -1$.\\
		Vậy tập xác định của hàm số là $(-1;+\infty)$.
	}
\end{ex}
%%%============EX_2==============%%%
\begin{ex}%[0-HK1-CD-1-SoBacNinh-2324]%[VN-MT-6, Tống Văn Ký]%[0D3H1-5]
	Cho hàm số $y=f(x)$ nghịch biến trên $(0;+\infty)$. Khẳng định nào sau đây đúng?
	\choice
	{$f(1) < f(2)$}
	{$f\left(\sqrt{3} \right) > f\left(\sqrt{2} \right)$}
	{$f\left(\sqrt{2} \right)\le f\left(\sqrt{5} \right)$}
	{\True $f\left(\sqrt{3} \right) > f\left(\sqrt{5} \right)$}
	\loigiai{
		Vì $f(x)$ nghịch biến trên khoảng $(0;+\infty )$ và $\sqrt{3} < \sqrt{5}$ nên $f\left(\sqrt{3} \right) > f\left(\sqrt{5} \right)$.
	}
\end{ex}
%%%============EX_3==============%%%
\begin{ex}%[0-HK1-CD-1-SoBacNinh-2324]%[VN-MT-6, Tống Văn Ký]%[0D3H1-4]
	Đồ thị hàm số $y=-x+3$ đi qua điểm nào sau đây?
	\choice
	{$M(-3;0)$}
	{\True $N(1;2)$}
	{$P(2;-1)$}
	{$Q(0;-3)$}
	\loigiai{
		Thay tọa độ các điểm vào hàm số, ta thấy đồ thị hàm số $y=-x+3$ đi qua điểm $N(1;2)$.
	}
\end{ex}
%%%============EX_4==============%%%
\begin{ex}%[0-HK1-CD-1-SoBacNinh-2324]%[VN-MT-6, Tống Văn Ký]%[0D3H1-3]
	Cho hàm số $f(x)=\heva{& -2x+1 & \text{ khi } x\le 0 \\
		& 2x^2-x+1 & \text{ khi } x > 0}$. Giá trị của $f(1)+f(-1)$ bằng
	\choice
	{$3$}
	{$2$}
	{\True $5$}
	{$0$}
	\loigiai{
		Ta có $f(-1)=-2\cdot (-1)+1=3$ và $f(1)=2-1+1=2$.\\
		Vậy $f(1)+f(-1)= 5$.
	}
\end{ex}
%%%============EX_5==============%%%
\begin{ex}%[0-HK1-CD-1-SoBacNinh-2324]%[VN-MT-6, Tống Văn Ký]%[0D3N2-1]
	Trong các hàm số sau, hàm số nào là hàm số bậc hai?
	\choice
	{$y=2x^2+5x+\dfrac{1}{x}$}
	{$y=-x^2+2\sqrt{x}+1$}
	{$y=-5+12x$}
	{\True $y=2x^2+x+2$}
	\loigiai{
		Hàm số $y=2x^2+x+2$ là hàm số bậc hai.
	}
\end{ex}
%%%============EX_6==============%%%
\begin{ex}%[0-HK1-CD-1-SoBacNinh-2324]%[VN-MT-6, Tống Văn Ký]%[0D3H2-3]
	Đỉnh của parabol $y=x^2-2x+5$ là điểm có tọa độ là
	\choice
	{\True $(1;4)$}
	{$(2;5)$}
	{$(-1;8)$}
	{$(0;5)$}
	\loigiai{
		Ta có $-\dfrac{b}{2a}=1$ và $f(1)=4$ nên đỉnh của parabol $y=x^2-2x+5$ là điểm có tọa độ là $(1;4)$.
	}
\end{ex}
%%%============EX_7==============%%%
\begin{ex}%[0-HK1-CD-1-SoBacNinh-2324]%[VN-MT-6, Tống Văn Ký]%[0D3H2-3]
	\immini[thm]{Cho hàm số bậc hai $y=ax^2+bx+c \left(a,b,c\in \mathbb{R}\right)$ có đồ thị như hình vẽ bên. Mệnh đề nào sau đây đúng?
		\choice
		{$a > 0, b > 0, c > 0$}
		{\True $a < 0, b < 0, c > 0$}
		{$a > 0, b < 0, c > 0$}
		{$a < 0, b < 0, c < 0$}
	}{
		\begin{tikzpicture}[scale=1, font=\footnotesize, line join=round, line cap=round, >=stealth]
			\def\xmin{-3}\def\xmax{1}
			\def\ymin{-0.5}\def\ymax{2.5}
			\draw[->] (\xmin,0)--(\xmax,0) node[below] {$x$};
			\draw[->] (0,\ymin)--(0,\ymax) node[right] {$y$};
			\draw (0,0) circle (1pt) node [above left] {$O$};
			\clip (\xmin,\ymin) rectangle (\xmax,\ymax);
			\draw[smooth,samples=200,domain=\xmin:\xmax] plot (\x,{-1*((\x)^2)-2*\x+1});		
		\end{tikzpicture}
		
	}
	\loigiai{
		Từ đồ thị, ta thấy
		\begin{itemize}
			\item Đồ thị có bề lõm hướng xuống dưới nên $a < 0$.
			\item Đồ thị cắt trục $Oy$ tại điểm có tung độ dương nên $c>0$.
			\item Đồ thị có đỉnh nằm bên trái $Oy$ nên $-\dfrac{b}{2a}<0\Rightarrow ab>0$. Mà $a<0$ nên $b<0$.
		\end{itemize}
	}
\end{ex}
%%%============EX_8==============%%%
\begin{ex}%[0-HK1-CD-1-SoBacNinh-2324]%[VN-MT-6, Tống Văn Ký]%[0D3H2-2]
	Hàm số $y=2x^2-4x-5$ đồng biến trên khoảng nào dưới đây?
	\choice
	{$(-\infty;-1)$}
	{$(-1;+\infty)$}
	{\True $(1;+\infty)$}
	{$(-\infty;1)$}
	\loigiai{
		Ta có $a=2>0$ và $-\dfrac{b}{2a}=1$ nên hàm số $y=2x^2-4x-5$ đồng biến trên khoảng $(1;+\infty)$.
	}
\end{ex}
%%%============EX_9==============%%%
\begin{ex}%[0-HK1-CD-1-SoBacNinh-2324]%[VN-MT-6, Tống Văn Ký]%[0D3H2-3]
	\immini[thm]{Cho hàm số bậc hai có đồ thị như hình vẽ bên. Hàm số đó là hàm số nào?
		\choice
		{$y=3x^2-2x$}
		{$y=-x^2-2x$}
		{\True $y=-x^2+2x$}
		{$y=-x^2+2x+1$}
	}{
		\begin{tikzpicture}[scale=1.2, font=\footnotesize, line join=round, line cap=round, >=stealth]
			\def\xmin{-1}\def\xmax{3}\def\ymin{-1}\def\ymax{1.5}
			\draw[->] (\xmin,0)--(\xmax,0) node[below] {$x$};
			\draw[->] (0,\ymin)--(0,\ymax) node[right] {$y$};
			\draw (0,0) circle (1pt) node [above left] {$O$};
			\clip (\xmin,\ymin) rectangle (\xmax,\ymax);
			\draw[smooth,samples=100,domain=\xmin:\xmax] plot (\x,{-1*((\x)^2)+2*\x+0});
			\draw[dashed] (1,0)--(1,1)--(0,1);
			\draw (0,1) node[left]{$1$};
			\draw (1,0) node[below]{$1$};		
		\end{tikzpicture}
	}
	\loigiai{
		Giả sử hàm số bậc hai cần tìm có dạng $y =ax^2+bx+c$.\\
		Vì đồ thị đi qua gốc tọa độ nên $c=0$.\\
		Vì đồ thị có đỉnh là $I(1; 1)$ nên 
		\[\heva{& -\dfrac{b}{2a}=1 \\& a+b=1} \Leftrightarrow \heva{& 2a+b =0 \\ & a+b=1} \Leftrightarrow \heva{& a=-1 \\& b=2.}\]
		Vậy $y=-x^2+2x$.
	}
\end{ex}
%%%============EX_10==============%%%
\begin{ex}%[0-HK1-CD-1-SoBacNinh-2324]%[VN-MT-6, Tống Văn Ký]%[0H5H3-2]
	\immini[thm]{Cho tam giác $ABC$ có $M$ là trung điểm $BC$. Mệnh đề nào sau đây \textbf{sai}?
		\choice[2]
		{\True $\overrightarrow{AB}-\overrightarrow{AC}=2\overrightarrow{BM}$}
		{$\overrightarrow{AB}+\overrightarrow{AC}=2\overrightarrow{AM}$}
		{$\overrightarrow{MB}+\overrightarrow{MC}=\overrightarrow0$}
		{$\overrightarrow{AC}-\overrightarrow{AB}=2\overrightarrow{BM}$}
	}{\begin{tikzpicture}[font=\footnotesize, line join=round, line cap=round, >=stealth, scale=1]
		\path 
			(0:0) coordinate (B)++(0:3) coordinate (C)
			($(B)!1/2!(C)$) coordinate (M)
			(B)++(60:3) coordinate (A)
			;
			\draw 
			(M)--(A)--(B)--(C)--(A)
			;
			\foreach \x/\g in {M/-90,C/0,A/90,B/180}
			\fill (\x) circle (1pt)
			+(\g:3mm) node{$\x$};
		\end{tikzpicture}
	}
	\loigiai{
		Ta có $\overrightarrow{AB}-\overrightarrow{AC} =\overrightarrow{CB}=2\overrightarrow{MB} =-2\overrightarrow{BM}$ nên khẳng định $\overrightarrow{AB}-\overrightarrow{AC} =2\overrightarrow{BM}$ là sai.
	}
\end{ex}
%%%============EX_11==============%%%
\begin{ex}%[0-HK1-CD-1-SoBacNinh-2324]%[VN-MT-6, Tống Văn Ký]%[0H5H3-2]
	\immini[thm]{Cho hình bình hành $ABCD$ tâm $I$. Khẳng định nào sau đây đúng?
		\choice[2]
		{$\overrightarrow{AB}+\overrightarrow{AC}+\overrightarrow{AD}=3\overrightarrow{AI}$}
		{\True $\overrightarrow{AB}+2\overrightarrow{AC}+\overrightarrow{AD}=3\overrightarrow{AC}$}
		{$\overrightarrow{IA}+\overrightarrow{IB}=\overrightarrow{IC}+\overrightarrow{ID}$}
		{$\overrightarrow{AB}=\overrightarrow{CD}$}
	}{\begin{tikzpicture}[font=\footnotesize, line join=round, line cap=round, >=stealth, scale=1]	
		\path 
			(0:0) coordinate (B)++(0:3) coordinate (C)
			(B)++(60:2.5) coordinate (A)
			($(A)+(C)-(B)$) coordinate (D)
			($(A)!1/2!(C)$) coordinate (I)
			;
			\draw 
			(D)--(A)--(B)--(C)--(D) (A)--(C) (B)--(D)
			;
			\foreach \x/\g in {D/60,C/0,A/90,B/180,I/-100}
			\fill (\x) circle (1pt)
			+(\g:3mm) node{$\x$};
		\end{tikzpicture}
	}
	\loigiai{
		Ta có $\overrightarrow{AB}+2\overrightarrow{AC}+\overrightarrow{AD} =\left(\overrightarrow{AB}+\overrightarrow{AD}\right)+2\overrightarrow{AC} =\overrightarrow{AC}+2\overrightarrow{AC} =3\overrightarrow{AC}$.
	}
\end{ex}
%%%============EX_12==============%%%
\begin{ex}%[0-HK1-CD-1-SoBacNinh-2324]%[VN-MT-6, Tống Văn Ký]%[0H5H3-6]
	Cho tam giác $ABC$ đều cạnh bằng $a$. Độ dài vectơ $\overrightarrow{AB}+\overrightarrow{AC}$ bằng
	\choice
	{$2a$}
	{$\dfrac{\sqrt{3} a}{2}$}
	{\True $a\sqrt{3}$}
	{$a\sqrt{2}$}
	\loigiai{
		\begin{center}
			\begin{tikzpicture}[font=\footnotesize, line join=round, line cap=round, >=stealth, scale=1]
				\path 
				(0:0) coordinate (B)++(0:3) coordinate (C)
				($(B)!1/2!(C)$) coordinate (M)
				(B)++(60:3) coordinate (A)
				;
				\draw 
				(M)--(A)--(B)--(C)--(A)
				;
				\foreach \x/\g in {M/-90,C/0,A/90,B/180}
				\fill (\x) circle (1pt)
				+(\g:3mm) node{$\x$};
			\end{tikzpicture}
		\end{center}
			Gọi $M$ là trung điểm $BC$.\\
			Ta có $\left|\overrightarrow{AB}+\overrightarrow{AC} \right| =\left|2\overrightarrow{AM}\right|=2AM$.\\
			Vì tam giác $ABC$ đều, có $M$ là trung điểm $BC$ nên $AM \perp BC$.\\
			Suy ra $AM=\sqrt{AB^2-BM^2}=\dfrac{a\sqrt{3}}{2}$.\\
			Vậy $\left|\overrightarrow{AB}+\overrightarrow{AC} \right|=2AM=a\sqrt{3}$.}
\end{ex}

\Closesolutionfile{ans}

%================================PHẦN II
\cauds
\Opensolutionfile{ans}[ans/0-HK1-CD-1-SoBacNinh-2324-Phan-II]
%==========Câu 1
\begin{ex}%[0-HK1-CD-1-SoBacNinh-2324]%[VN-MT-6, Tống Văn Ký]%[0D3H1-2]
	Cho hàm số $y=\sqrt{3-2x}$ có tập xác định là $\mathscr{D}$.
	\choiceTF
	{$\mathscr{D}=\left[\dfrac{3}{2};+\infty\right)$}
	{\True $x_0=1\in \mathscr{D}$}
	{\True $(-3;0) \subset \mathscr{D}$}
	{$\left(\dfrac{3}{2};3\right) \cap \mathscr{D}=\left\{\dfrac{3}{2}\right\}$}
	\loigiai{
		\begin{itemchoice}
		\itemch {\bf Sai.}\\ 
		Điều kiện xác định là $3-2x \geq 0\Leftrightarrow x \leq \dfrac{3}{2}$. Suy ra $\mathscr{D}=\left(-\infty;\dfrac{3}{2}\right]$.
		\itemch {\bf Đúng.}\\ 
		Ta có $\mathscr{D}=\left(-\infty; \dfrac{3}{2}\right]$ nên $x_0=1\in \mathscr{D}$.
		\itemch {\bf Đúng.}\\ 
		Vì $(-3;0) \subset\left(-\infty;\dfrac{3}{2}\right]$ nên $(-3;0) \subset\mathscr{D}$.
		\itemch {\bf Sai.}\\ 
		Ta có $\left(\dfrac{3}{2};3\right) \cap \mathscr{D}=\varnothing$.
		\end{itemchoice}
	}
\end{ex}
\begin{ex}%[0-HK1-CD-1-SoBacNinh-2324]%[VN-MT-6, Tống Văn Ký]%[0H5H3-2]
	Cho tam giác $ABC$ có $G$ là trọng tâm. Gọi $G'$ là điểm đối xứng với $G$ qua trung điểm $M$ của $BC$. Khi đó
	\choiceTF
	{\True Hai vectơ $\overrightarrow{GA}$, $\overrightarrow{G'G}$ cùng hướng}
	{$\overrightarrow{BM}=\overrightarrow{CM}$}
	{\True $\overrightarrow{GA}=\overrightarrow{G'G}$}
	{\True $\overrightarrow{BG'}=\overrightarrow{GC}$}
	\loigiai{
		\begin{center}
			\begin{tikzpicture}[font=\footnotesize, line join=round, line cap=round, >=stealth, scale=1, declare function={a=5; goc=70; b=0.65*a;}]
				\path
				(goc:b) coordinate (A)
				(0,0) coordinate (B)
				(a,0) coordinate (C)
				($(A)!0.5!(B)$) coordinate (D)
				($(B)!0.5!(C)$) coordinate (M)
				($(C)!0.5!(A)$) coordinate (F)
				(intersection of A--M and B--F) coordinate (G)
				($(G)!2!(M)$) coordinate (G')
				;
				\foreach \x/\y in {G/M,M/G'}{
					\path (\x)--(\y) node[midway,sloped]{\tikz{\draw (-90:2pt)--(90:2pt);}};
				}
				\draw 
				(A)--(B)--(C)--cycle
				(A)--(G')
				;
				\foreach \t/\g in {A/90,B/-120,C/-60,G/60,M/-120,G'/0}{
					\draw[fill=white] (\t) circle (1pt) node[shift={(\g:7pt)},font=\scriptsize]{$ \t $};
				}
			\end{tikzpicture}	
		\end{center}
		\begin{itemchoice}
			\itemch {\bf Đúng.}\\ Dễ thấy hai vectơ $\overrightarrow{GA}$, $\overrightarrow{G'G}$ cùng hướng.
			\itemch {\bf Sai.}\\ $\overrightarrow{BM}$, $\overrightarrow{CM}$ ngược hướng nên không bằng nhau. 
			\itemch {\bf Đúng.}\\ Vì $M$ là trung điểm của $BC$, $G$ là trọng tâm tam giác $ABC$ nên $\overrightarrow{GA}=2\overrightarrow{MG}$.\\
			Lại có $G'$ đối xứng với $G$ qua $M$ nên $\overrightarrow{G'G}=2\overrightarrow{MG}$.\\
			Vậy $\overrightarrow{GA}=\overrightarrow{G'G}$.
			\itemch {\bf Đúng.}\\ 
			Tứ giác $BGCG'$ là hình bình hành do hai đường chéo cắt nhau tại trung điểm mỗi đường nên $\overrightarrow{BG'}=\overrightarrow{GC}$.	
		\end{itemchoice}
	} 
\end{ex}
%==========Câu 4
\begin{ex}%[0-HK1-CD-1-SoBacNinh-2324]%[VN-MT-6, Tống Văn Ký]%[0D3H2-4]
	Cho hàm số $y=x^2-2x-3$ có đồ thị là parabol $(P)$.
	\choiceTF
	{\True Parabol $(P)$ có đỉnh là $I(1;-4)$}
	{Parabol $(P)$ có bề lõm hướng xuống dưới}
	{\True Parabol $(P)$ cắt trục hoành tại hai điểm phân biệt}
	{\True Hàm số $y=x^2-2x-3$ nghịch biến trên $(-\infty;1)$ và đồng biến trên $(1;+\infty)$}
	\loigiai{
		\begin{itemchoice}
		\itemch {\bf Đúng.}\\ Ta có $-\dfrac{b}{2a}=1$ suy ra $(P)$ có đỉnh là $I(1;-4)$.
		\itemch {\bf Sai.}\\ Hệ số $a=1> 0$, nên parabol có bề lõm hướng lên trên.
		\itemch {\bf Đúng.}\\ Ta có $(P)$ giao $Ox$ tại các điểm $(-1;0)$ và $(3;0)$.
		\itemch {\bf Đúng.}\\ Ta có $-\dfrac{b}{2a}=1$ nên hàm số nghịch biến trên $(-\infty;1)$ và đồng biến trên $(1;+\infty)$.
			\end{itemchoice}
	}
\end{ex}
\begin{ex}%[0-HK1-CD-1-SoBacNinh-2324]%[VN-MT-6, Tống Văn Ký]%[0H5V2-6]
	Cho ba lực $\overrightarrow{F}_1=\overrightarrow{MA}$, $\overrightarrow{F}_2=\overrightarrow{MB}$, $\overrightarrow{F}_3=\overrightarrow{MC}$ cùng tác động vào một vật tại điểm $M$ và vật đứng yên như hình vẽ.
	\begin{center}
		\begin{tikzpicture}[font=\footnotesize, line join=round, line cap=round, >=stealth, scale=1]
			\def\a{2}
			\pgfmathsetmacro\b{\a*sqrt(3)}
			\path[draw] (0,0)coordinate(M)
			(60:\a)coordinate(B)
			(-30:\b)coordinate(C)
			($(B)+(C)$)coordinate(D)
			($(M)-(D)$) coordinate(A)
			; 
			\draw[->] (M)--(B);
			\draw[->] (M)--(A);
			\draw[->] (M)--(C);		
			\foreach \d/\g in {A/90,B/180,C/0,M/0}\fill(\d) circle (1pt)+(\g:.3)node{$\d$};
			\begin{scope}
				\clip (B)--(M)--(A)--cycle;
				\draw (M) circle(3mm);
			\end{scope}
			\node at ($(M)+(160:7mm)$) {$120^{\circ}$};
			\begin{scope}
				\clip (A)--(M)--(C)--cycle;
				\draw (M) circle(4mm);
			\end{scope}
			\node at ($(M)+(260:7mm)$) {$150^{\circ}$};
		\end{tikzpicture}
	\end{center}
	Biết cường độ của lực $\overrightarrow{F}_1$ là $50$ N, $\widehat{AMB}=120^{\circ}$, $\widehat{AMC}=150^{\circ}$. Cường độ của lực $\overrightarrow{F}_3$ là	
	\choiceTF
	{\True $\left|\overrightarrow{F}_3\right|=25\sqrt{3}$ N}
	{\True $\overrightarrow{F}_1+\overrightarrow{F}_2+\overrightarrow{F}_3=\overrightarrow{0}$}
	{$\left|\overrightarrow{F}_2\right|=25\sqrt{3}$ N}
	{\True $\left|\overrightarrow{F}_1\right|=50$ N}
	\loigiai
	{\begin{center}
			\begin{tikzpicture}[font=\footnotesize, line join=round, line cap=round, >=stealth, scale=1]
				\def\a{2}
				\pgfmathsetmacro\b{\a*sqrt(3)}
				\path[draw] (0,0)coordinate(M)
				(60:\a)coordinate(B)
				(-30:\b)coordinate(C)
				($(B)+(C)$)coordinate(D)
				($(M)-(D)$) coordinate(A)
				; 
				\draw[->] (M)--(B);
				\draw[->] (M)--(A);
				\draw[->] (M)--(C);
				\draw[->] (M)--(D);
				\draw (B)--(D)--(C);				
				\foreach \d/\g in {A/90,B/180,C/0,M/30,D/0}\fill(\d) circle (1pt)+(\g:.3)node{$\d$};
				\begin{scope}
					\clip (B)--(M)--(A)--cycle;
					\draw (M) circle(3mm);
				\end{scope}
				\node at ($(M)+(160:7mm)$) {$120^{\circ}$};
				\begin{scope}
					\clip (A)--(M)--(C)--cycle;
					\draw (M) circle(4mm);
				\end{scope}
				\node at ($(M)+(260:7mm)$) {$150^{\circ}$};
			\end{tikzpicture}
		\end{center}
		Ta có $\widehat{AMB}=120^{\circ}$, $\widehat{AMC}=150^{\circ}$ $\Rightarrow \widehat{BMC}=360^{\circ}-120^{\circ}-150^{\circ}=90^{\circ}$.\\
		Vẽ hình chữ nhật $MCDB$ có $\widehat{CMD}=180^{\circ}-\widehat{AMC}=180^{\circ}-150^{\circ}=30^{\circ}$.\\
		Vì vật đứng yên nên tổng hợp lực tác động vào vật bằng $\overrightarrow{0}$.
		\begin{itemchoice}
		\itemch {\bf Đúng.}\\ Ta có $\cos \widehat{CMD}=\dfrac{MC}{MD} \Rightarrow MC=MD \cdot \cos 30 =50 \cdot \dfrac{\sqrt{3}}{2}=25\sqrt{3}$.\\
			Vậy $\left|\overrightarrow{F}_3\right|=MC=25\sqrt{3}$ N.
		\itemch {\bf Đúng.}\\ Do vật đứng yên nên $\overrightarrow{F}_1+\overrightarrow{F}_2+\overrightarrow{F}_3=\overrightarrow{0}$.
		\itemch {\bf Sai.}\\ Lại có $\sin \widehat{CMD}=\dfrac{CD}{MD} \Rightarrow MB=CD=MD \cdot \sin 30^{\circ}=50 \cdot \dfrac{1}{2}=25$.\\
			Vậy $\left|\overrightarrow{F}_2\right|=MB=25$ N.
		\itemch {\bf Đúng.}\\ Theo đề cho, ta có $\left|\overrightarrow{F}_1\right|=50$ N.
		\end{itemchoice}
	}
\end{ex}
\Closesolutionfile{ans}

%================================PHẦN III
\caukq
\Opensolutionfile{ans}[ans/0-HK1-CD-1-SoBacNinh-2324-Phan-III]
%%% CÂU 1
\begin{ex}%[0-HK1-CD-1-SoBacNinh-2324]%[VN-MT-6, Tống Văn Ký]%[0D3N1-3]
	Cho hàm số $y=\heva{&\dfrac{2}{x-1}, &x \in(-\infty; 0) \\
		&\sqrt{x+1}, &x \in[0; \infty)}$. Giá trị $f(0)+f(-1)$ bằng bao nhiêu?
	\shortans[]{$0$} %Chuẩn hóa ko quá 4 ký tự
\loigiai{
	 Với $x=0$ ta có $f(0)=\sqrt{0+1}=1$.\\
	Với $x=-1$ ta có $f(-1)=-1$.\\
	Vậy $f(0)+f(-1)=0$.}
\end{ex}
\begin{ex}%[0D3N1-5]
	Cho hàm số $y=(m-7)x+2$ ($m$ là tham số thực). Có bao nhiêu giá trị nguyên dương của tham số $m$ để hàm số đã cho là hàm số nghịch biến trên $\mathbb{R}$.
		\shortans[]{$6$} %Chuẩn hóa ko quá 4 ký tự
	\loigiai{
		Hàm số đã cho nghịch biến trên $\mathbb{R}$ khi và chỉ khi $m-7<0\Leftrightarrow m<7$.\\
		Mà $m$ là số nguyên dương nên $m \in\{1;2;3;4;5;6\}$.\\
		Vậy có $6$ giá trị nguyên dương của tham số $m$ để hàm số đã cho là hàm số nghịch biến trên $\mathbb{R}$.}
\end{ex}
\begin{ex}%[0-HK1-CD-1-SoBacNinh-2324]%[VN-MT-6, Tống Văn Ký]%[0D3H1-2]
	Biết rằng tập xác định của hàm số $ y=\sqrt{x-2}+\dfrac{3}{\sqrt{5-x}}$ là $\mathscr{D} =[a; b)$. Tính giá trị của $b -a$.
	\shortans[]{$3$} %Chuẩn hóa ko quá 4 ký tự
\loigiai{
		Hàm số $y=\sqrt{x-2}+\dfrac{3}{\sqrt{5-x}}$ xác định $\Leftrightarrow \heva{ & x-2\ge 0 \\ & 5-x > 0 } \Leftrightarrow \heva{ & x\ge 2 \\ & x < 5 }\Leftrightarrow 2 \le x < 5$.\\
		Tập xác định của hàm số là $\mathscr{D}=[2;5)$. Suy ra $a=2$, $b=5$ và $b-a=3$.}
\end{ex}

\begin{ex}%[0-HK1-CD-1-SoBacNinh-2324]%[VN-MT-6, Tống Văn Ký]%[0D3H2-3]
	 Cho hàm số $y=ax^2+bx+2$ có đồ thị là parabol $(P)$. Tính giá trị của biểu thức $S=a+b$, biết $(P)$ có đỉnh là điểm $I(2;6)$.
	\shortans[]{$3$} %Chuẩn hóa ko quá 4 ký tự
\loigiai{
	 Parabol $y=ax^2+bx+2$ có đỉnh là $I(2;6)$, nên ta có hệ phương trình
			\[\heva{ & -\dfrac{b}{2a}=2 \\ & y(2)=6} \Leftrightarrow \heva{ & 4a+b=0 \\ & 4a+2b+2=6} \Leftrightarrow \heva{& a=-1 \\ & b=4.}
			\]
			Vậy hàm số là $y=-x^2+4x+2$. Do đó $S=a+b=3$.}
\end{ex}

\begin{ex}%[0-HK1-CD-1-SoBacNinh-2324]%[VN-MT-6, Tống Văn Ký]%[0H5V3-8]
	Cho tam giác $ABC$ vuông tại $A$, có $AB=3, BC=5$. Gọi $E$ là trung điểm $AB$. Gọi $M$ là điểm thay đổi trên đường thẳng $BC$. Tìm giá trị nhỏ nhất của biểu thức 
	\[T=\left|\overrightarrow{MA}+\overrightarrow{MB}+2\overrightarrow{MC}\right|.\]
	\shortans[]{$2{,}4$} %Chuẩn hóa ko quá 4 ký tự
\loigiai{
		\begin{center}
			\begin{tikzpicture}[font=\footnotesize, line join=round, line cap=round, >=stealth, scale=1]
				\path 
				(0:0) coordinate (A)++(0:4) coordinate (C)
				(A)++(90:3) coordinate (B)
				($(B)!1/2!(A)$) coordinate (E)
				($(E)!1/2!(C)$) coordinate (I)
				($(B)!9/25!(C)$) coordinate (H)
				($(B)!1/2!(H)$) coordinate (J)
				($(J)!1/2!(C)$) coordinate (K)
				;
				\draw 
				(A)--(B)--(C)--(A) (C)--(E) (E)--(J) (A)--(H) (I)--(K)
				;
				\pic[draw,angle radius=2mm] {right angle=E--J--B}
				pic[draw,angle radius=2mm] {right angle=A--H--B}
				pic[draw,angle radius=2mm] {right angle=I--K--B}
				;
				\foreach \x/\g in {E/180,B/100,A/-150,C/-20,I/-120,J/60, H/60, K/60}
				\fill (\x) circle (1pt)
				+(\g:3mm) node{$\x$};
			\end{tikzpicture}
		\end{center}
			Vì $E$ là trung điểm $AB$ nên $\overrightarrow{AB}=2\overrightarrow{AE}$.\\
			Suy ra $2\overrightarrow{AE}+\overrightarrow{CA}=\overrightarrow{AB}+\overrightarrow{CA}=\overrightarrow{CB}$.\\	
			Gọi $I$ là điểm thỏa mãn $\overrightarrow{IA}+\overrightarrow{IB}+2\overrightarrow{IC}=\overrightarrow0\Leftrightarrow 2\overrightarrow{IE}+2\overrightarrow{IC}=\overrightarrow0\Leftrightarrow \overrightarrow{IE}+\overrightarrow{IC}=\overrightarrow0$.\\	
			Suy ra $I$ là trung điểm của $EC$.
			Ta có
			\begin{eqnarray*}
				T &=&\left|\overrightarrow{MA}+\overrightarrow{MB}+2\overrightarrow{MC}\right|=\left|\overrightarrow{MI}+\overrightarrow{IA}+\overrightarrow{MI}+\overrightarrow{IB}+2\left(\overrightarrow{MI}+\overrightarrow{IC}\right)\right| \\
				&=& \left|4\overrightarrow{MI}+\overrightarrow{IA}+\overrightarrow{IB}+2\overrightarrow{IC}\right|\\
				&=& \left|4\overrightarrow{MI}+\overrightarrow{0}\right|=\left|4\overrightarrow{MI}\right|=4MI.
			\end{eqnarray*}
			Tam giác $ABC$ vuông tại $A$ nên $AC=\sqrt{BC^2-AB^2}=4$.\\
			Do $M$ di động trên đường thẳng $BC$, nên $T$ đạt giá trị nhỏ nhất khi $M\equiv K$ là hình chiếu của $I$ trên đường thẳng $BC$.\\	
			Gọi $J$, $H$ lần lượt là hình chiếu vuông góc của $E$ và $A$ xuống cạnh $BC$.\\
			Suy ra $IK$, $AH$, $EJ$ đôi một song song với nhau.\\
			Dễ thấy $IK=\dfrac{1}{2} EJ=\dfrac{1}{4} AH$.\\	
			Vậy $\min T=4IK=AH=\dfrac{AB\cdot AC}{BC}=\dfrac{3\cdot 4}{5}=\dfrac{12}{5}=2{,}4$.
		}
\end{ex}

\begin{ex}%[0-HK1-CD-1-SoBacNinh-2324]%[VN-MT-6, Tống Văn Ký]%[0D3V2-6]
	Bác Quân là chủ một xưởng cơ khí sản xuất các thiết bị máy. Trong xưởng của bác, mỗi sản phẩm A có chi phí sản xuất là $10$ triệu đồng và dự định bán ra với giá là $15$ triệu đồng. Với giá bán đó, số sản phẩm mà bên khách hàng đối tác sẽ mua trong một năm là $600$ sản phẩm. Nhằm mục tiêu đẩy mạnh sản xuất và tiêu thụ sản phẩm này, bác Quân dự định giảm giá bán và ước tính rằng nếu cứ giảm $1$ triệu đồng mỗi sản phẩm thì số lượng sản phẩm bán ra trong một năm là sẽ tăng thêm $200$ sản phẩm. Bác Quân phải định giá bán mới của sản phẩm A là bao nhiêu, để sau khi đã thực hiện giảm giá, lợi nhuận thu được sẽ là cao nhất?
	\shortans[]{$14$} %Chuẩn hóa ko quá 4 ký tự
\loigiai{
			Gọi $x$ (triệu đồng) là số tiền dự định giảm giá của sản phẩm $(0\le x\le 5)$. Khi đó\\	
			Lợi nhuận thu được khi bán mỗi sản phẩm là $15-x-10$ $=5-x$ (triệu đồng).\\	
			Số sản phẩm mà xưởng sẽ bán được trong một năm là $600+200x$ (sản phẩm).\\	
			Lợi nhuận (đơn vị: triệu đồng) mà xưởng thu được trong một năm là
			\[ f(x)=(5-x)(600+200x) =-200x^2+400x+3\,000. \]
			Xét hàm số $f(x)=-200x^2+400x+3\,000=-200(x-1)^2+3\,200\le 3\,200$.\\	
			Dấu \lq\lq $=$\rq\rq\, xảy ra khi $x=1\in [0;5]$ nên $f(x)$ đạt giá trị lớn nhất bằng $3\,200$ khi $x=1$.\\	
			Vậy giá mới của mỗi sản phẩm A là $14$ triệu đồng thì lợi nhuận thu được là cao nhất.
	}
\end{ex}

\Closesolutionfile{ans}
% \newpage
% \begin{center}
% 	\fbox{\textbf{\color{purple} \large BẢNG ĐÁP ÁN}}
% \end{center}

% \noindent\textbf{PHẦN I.}
% \inputansbox[2]{6}{ans/0-HK1-CD-1-SoBacNinh-2324-Phan-I}

% \noindent\textbf{PHẦN II.}
% \inputansbox[2]{2}{ans/0-HK1-CD-1-SoBacNinh-2324-Phan-II}

% \noindent\textbf{PHẦN III.}
% \inputansbox[2]{6}{ans/0-HK1-CD-1-SoBacNinh-2324-Phan-III}


